
\documentclass{article}
\usepackage{colortbl}
\usepackage{makecell}
\usepackage{multirow}
\usepackage{supertabular}

\begin{document}

\newcounter{utterance}

\centering \large Interaction Transcript for game `codenames', experiment `concrete', episode 3 with reflexion{-}Qwen3.5{-}9B.
\vspace{24pt}

{ \footnotesize  \setcounter{utterance}{1}
\setlength{\tabcolsep}{0pt}
\begin{supertabular}{c@{$\;$}|p{.15\linewidth}@{}p{.15\linewidth}p{.15\linewidth}p{.15\linewidth}p{.15\linewidth}p{.15\linewidth}}
    \# & $\;$A & \multicolumn{4}{c}{Game Master} & $\;\:$B\\
    \hline

    \theutterance \stepcounter{utterance}  
    & & & \multicolumn{2}{p{0.3\linewidth}}{
        \cellcolor[rgb]{0.9,0.9,0.9}{
            \makecell[{{p{\linewidth}}}]{
                \texttt{\tiny{[GM$|$GM]}}
                {'hidden': {'team': ['marbles', 'clay'], 'innocent': ['armchair', 'cauliflower', 'underwear'], 'opponent': [], 'assassin': []}, 'revealed': {'team': {'team': [], 'innocent': [], 'opponent': [], 'assassin': []}, 'opponent': {'team': [], 'innocent': [], 'opponent': [], 'assassin': []}}}
            }
        }
    }
    & & \\ \\

    \theutterance \stepcounter{utterance}  
    & & \multicolumn{4}{p{0.6\linewidth}}{
        \cellcolor[rgb]{0.9,0.9,0.9}{
            \makecell[{{p{\linewidth}}}]{
                \texttt{\tiny{[P1$\langle$GM]}}
                \texttt{You have attempted this same task before and did not succeed. These are the plans you wrote after those attempts.} \\
\texttt{Trial \#1: New plan: The previous attempt failed because the clue "food" was too broad and included "marbles" and "clay," which are not edible, likely confusing the partner or wasting turns. Additionally, "vodka" is a drink, not strictly "food," making the association weak. Next time, I will choose a more precise semantic link that covers all selected targets without ambiguity. For instance, using "animal" for "donkey," "warthog," and "birdseed" (as bird food), or "container" for "paperclip" and "clay" (if shaped), but a better strategy is to group words by a stronger, undeniable category like "things that roll" (marbles, donkey, warthog) or "items made of natural materials" (clay, birdseed, donkey). I will ensure every target word in the list has a direct, unambiguous relationship with the clue word to maximize the partner's success rate.} \\
\\ 
\texttt{Trial \#2: New plan: The previous attempt failed because the clue "round" was ambiguous; while it fits "marbles" and "clay," it does not fit "paperclip" (which is flat), yet the partner likely guessed "paperclip" anyway, wasting a turn or triggering an opponent word. Additionally, grouping "marbles" and "clay" together is inefficient as they are distinct objects. Next time, I will choose a clue that is a perfect, unambiguous super{-}category for the specific subset of remaining words, avoiding broad terms like "round" or "small" that might include non{-}targets or fail to include the target. For the final two words "marbles" and "clay," I will use a more specific link like "spheres" (if clay can be molded into spheres) or, better yet, group them with a third word if possible, or simply pick one high{-}probability word with a very specific clue (e.g., "game" for marbles) to ensure a quick win without risking ambiguity. Given only two words left, I will prioritize a clue that strictly defines both, such as "moldable" (clay) and "game" (marbles) are not shared, so I must pick a clue that covers both: perhaps "round objects" is still risky. Actually, "marbles" and "clay" share the property of being "round" only if clay is shaped. A better strategy is to guess "marbles" with clue "game" and "clay" with clue "moldable" in separate turns if needed, but since I must group them, I will try "art supplies" (clay) or "toys" (marbles). Wait, they don't share a strong category. I will switch strategy: instead of grouping the last two difficult words, I will have guessed one earlier or chosen a clue that definitely covers both, like "round" but acknowledging the risk. Actually, the best fix is to avoid grouping words that don't share a strong semantic link. If I must group "marbles" and "clay", I will use "round" but accept the risk, or better, realize that "clay" can be "molded" and "marbles" are "rolled". No common word. I will change the plan to: always select a clue that is a direct definition of the group, e.g., for "marbles" and "clay", use "round" but ensure the partner knows "clay" implies a ball. If that fails, I should have grouped "marbles" with "clay" under "round" but the partner might have guessed "hubcap" (distractor) or "ball" (not in list). The real mistake was grouping "paperclip" with "marbles" and "clay" under "small" or "metal" earlier} \\
\\ 
\texttt{Trial \#3: New plan: The previous attempt failed because the clue "round" was ambiguous; while it fits "marbles," it does not semantically define "clay" (which is a material, not a shape), likely causing the partner to hesitate, guess a distractor, or trigger an opponent word. Grouping two words with a weak or conditional link (clay only being round if molded) is inefficient and risky. Next time, I will avoid grouping words that lack a direct, unambiguous super{-}category. Instead, I will select a single target word with a very specific, high{-}probability clue (e.g., "game" for marbles or "pottery" for clay) to ensure an immediate guess without risking ambiguity or wasting turns on a partner's confusion.} \\
\\ 
\\ 
\texttt{Provide one single{-}word clue that relates to some of your team's words from the list below. You can choose to how many and to which words your clue relates to.} \\
\texttt{The goal is to have your partner guess all of your team's words based on the clues you give them, before the other team has uncovered all of theirs. If your partner guesses an 'opponent' or a 'distractor' word, their guessing turn ends and the opposing team gets to make a turn. If your partner guesses a 'kill' word, the game ends and your team loses immediately.} \\
\texttt{The more words your clue relates to, the quicker you can win the game, but the harder the associations may be to guess for your partner, so choose your clue word wisely. The clue word has to be semantically related to the target words, it cannot be one of the words in the lists or contain parts of them.} \\
\\ 
\texttt{Always give your single{-}word clue and your comma{-}separated list of related target words in the following format and make your answers as short as possible, never include any other text than is required in this form:} \\
\\ 
\texttt{CLUE: <WORD>} \\
\texttt{TARGETS: <TARGETS>} \\
\\ 
\texttt{Your team words are: marbles, clay.} \\
\texttt{Your opponent's team words are: .} \\
\texttt{Distractor words are: armchair, cauliflower, underwear.} \\
\texttt{The kill words are: .} \\
            }
        }
    }
    & \\ \\

    \theutterance \stepcounter{utterance}  
    & \multicolumn{4}{p{0.6\linewidth}}{
        \cellcolor[rgb]{0.9,0.9,0.9}{
            \makecell[{{p{\linewidth}}}]{
                \texttt{\tiny{[P1$\rangle$GM]}}
                \texttt{CLUE: round} \\
\texttt{TARGETS: marbles, clay} \\
            }
        }
    }
    & & \\ \\

    \theutterance \stepcounter{utterance}  
    & & & \multicolumn{2}{p{0.3\linewidth}}{
        \cellcolor[rgb]{0.9,0.9,0.9}{
            \makecell[{{p{\linewidth}}}]{
                \texttt{\tiny{[GM$|$GM]}}
                \texttt{round} \\
            }
        }
    }
    & & \\ \\

    \theutterance \stepcounter{utterance}  
    & & & \multicolumn{2}{p{0.3\linewidth}}{
        \cellcolor[rgb]{0.9,0.9,0.9}{
            \makecell[{{p{\linewidth}}}]{
                \texttt{\tiny{[GM$|$GM]}}
                ['marbles', 'clay']
            }
        }
    }
    & & \\ \\

    \theutterance \stepcounter{utterance}  
    & & & \multicolumn{2}{p{0.3\linewidth}}{
        \cellcolor[rgb]{0.9,0.9,0.9}{
            \makecell[{{p{\linewidth}}}]{
                \texttt{\tiny{[GM$|$GM]}}
                {'word': 'marbles', 'assignment': 'team'}
            }
        }
    }
    & & \\ \\

    \theutterance \stepcounter{utterance}  
    & & & \multicolumn{2}{p{0.3\linewidth}}{
        \cellcolor[rgb]{0.9,0.9,0.9}{
            \makecell[{{p{\linewidth}}}]{
                \texttt{\tiny{[GM$|$GM]}}
                {'word': 'clay', 'assignment': 'team'}
            }
        }
    }
    & & \\ \\

    \theutterance \stepcounter{utterance}  
    & & & \multicolumn{2}{p{0.3\linewidth}}{
        \cellcolor[rgb]{0.9,0.9,0.9}{
            \makecell[{{p{\linewidth}}}]{
                \texttt{\tiny{[GM$|$GM]}}
                \texttt{opponent has won} \\
            }
        }
    }
    & & \\ \\

    \theutterance \stepcounter{utterance}  
    & \multicolumn{4}{p{0.6\linewidth}}{
        \cellcolor[rgb]{0.95,0.95,0.95}{
            \makecell[{{p{\linewidth}}}]{
                \texttt{\tiny{[P2$\rightarrow$P1]}}
                \texttt{Post{-}episode Reflexion (not sent to GM):} \\
\texttt{New plan: The previous attempt failed because the clue "round" was semantically weak for "clay" (a material) and potentially ambiguous, leading to a guess error or wasted turn. Grouping "marbles" and "clay" under a single vague category was inefficient. Next time, I will abandon grouping difficult pairs and instead issue a single, highly specific clue for one remaining word to guarantee an immediate correct guess. For "marbles," I will use the clue "game" (or "toy"), and for "clay," I will use "pottery" (or "moldable"), ensuring the clue is a direct, unambiguous definition of the target word without relying on conditional properties or broad categories.} \\
            }
        }
    }
    & & \\ \\

\end{supertabular}
}

\end{document}
